\Needspace{9\baselineskip}\begin{center}\textit{Zwei Verfahren zur Unterscheidung von zusammengesetzten Zahlen und Primzahlen und zur Ermittlung der Faktoren der ersteren.}\end{center}


\Needspace{6\baselineskip}\Article{329}
Dass die Aufgabe, Primzahlen von zusammengesetzten Zahlen zu unterscheiden und die letzteren in ihre Primfaktoren zu zerlegen, zu den wichtigsten und nützlichsten der gesamten Arithmetik gehört und den Fleiß und Scharfsinn der Mathematiker sowohl der älteren als auch der neueren Zeit beschäftigt hat, ist so bekannt, dass es überflüssig wäre, ausführlich darüber zu sprechen. Dennoch müssen wir zugeben, dass sämtliche bisher vorgeschlagenen Verfahren entweder auf sehr besondere Fälle beschränkt oder so mühsam und langwierig sind, dass sie schon bei Zahlen innerhalb der Grenzen der von verdienstvollen Männern aufgestellten Tafeln, also bei Zahlen, für die aufwendige Verfahren überflüssig sind, selbst die Geduld eines geübten Rechners erschöpfen; auf größere Zahlen lassen sie sich meist kaum anwenden. Jene allgemein verfügbaren Tafeln, von denen wir hoffen dürfen, dass sie nach und nach weiter fortgesetzt werden, reichen zwar in den meisten gewöhnlichen Fällen gewiss aus. Dennoch bietet sich einem erfahrenen Rechner nicht selten Gelegenheit, aus der Zerlegung großer Zahlen in Faktoren erhebliche Vorteile zu ziehen, die einen mäßigen Zeitaufwand reichlich aufwiegen. Außerdem scheint es die Würde der Wissenschaft zu verlangen, sämtliche Hilfsmittel zur Lösung eines so eleganten und berühmten Problems sorgfältig weiterzuentwickeln. Aus diesen Gründen
\SourceLeaf{407}
zweifeln wir nicht, dass die beiden folgenden Verfahren, deren Wirksamkeit und Kürze wir aus langer Erfahrung bestätigen können, den Liebhabern der Arithmetik willkommen sein werden. Im Übrigen liegt es in der Natur der Aufgabe, dass \textit{jedes} Verfahren zunehmend aufwendiger wird, je größer die Zahlen sind, auf die es angewandt wird. Bei den folgenden Verfahren wachsen die Schwierigkeiten jedoch sehr langsam; Zahlen mit sieben, acht oder sogar noch mehr Ziffern sind besonders durch das zweite stets mit Erfolg behandelt worden, und zwar mit der ganzen Geschwindigkeit, die man bei so großen Zahlen vernünftigerweise erwarten darf. Nach sämtlichen bisher bekannten Verfahren würden sie eine Arbeit verlangen, die selbst einem unermüdlichen Rechner unerträglich wäre.

Bevor man die folgenden Verfahren verwendet, ist es stets sehr nützlich, die Division der gegebenen Zahl durch einige der kleinsten Primzahlen zu versuchen, etwa durch \(2,\allowbreak 3,\allowbreak 5,\allowbreak 7\) usw. bis \(19\) oder noch weiter. So vermeidet man nicht nur, es zu bereuen, einen solchen Teiler durch schwierige und aufwendige Verfahren gefunden zu haben, obwohl er sich durch bloße Division viel leichter hätte finden lassen\HistoricalNote{*)}{Umso mehr, als sich im Allgemeinen unter sechs Zahlen kaum eine findet, die durch keine der Zahlen \(2,\allowbreak 3,\allowbreak 5\ldots19\) teilbar ist.}; auch kann das zweite Verfahren, wenn keine dieser Divisionen gelingt, die dabei entstehenden \textit{Reste} mit großem Nutzen verwenden. Ist zum Beispiel die Zahl \(314159265\) in ihre Faktoren zu zerlegen, so gelingt die Division durch \(3\) zweimal, danach auch die Division durch \(5\) und \(7\), sodass man \(314159265=9.5.7.997331\) erhält. Daher genügt es, die Zahl \(997331\), die durch \(11,\allowbreak 13,\allowbreak 17,\allowbreak 19\) nicht teilbar ist, genauer zu untersuchen. Ebenso werden wir bei der gegebenen Zahl \(43429448\) den Faktor \(8\) abtrennen und die aufwendigeren Verfahren auf den Quotienten \(5428681\) anwenden.

\Needspace{6\baselineskip}\Article{330}
Das \textsc{erste Verfahren} beruht auf dem Satz, dass \textit{jede positive oder} \textit{negative Zahl, die quadratischer Rest einer anderen Zahl} \(M\) \textit{ist, auch ein Rest jedes} \textit{Teilers von} \(M\) \textit{ist.} Bekanntlich ist, wenn \(M\) durch keine Primzahl unterhalb \(\surd M\) teilbar ist, \(M\) gewiss eine Primzahl. Sind aber sämtliche Primzahlen unterhalb dieser Grenze, die \(M\) teilen, gleich \(p,\allowbreak q\) usw., so besteht \(M\) entweder aus diesen \textit{allein} oder ihren Potenzen oder kann nur \textit{einen} weiteren Primfaktor größer als \(\surd M\) enthalten. Diesen findet man, indem man \(M\) durch \(p,\allowbreak q\) usw. so oft wie möglich dividiert. Bezeichnen wir die Gesamtheit aller Primzahlen unterhalb \(\surd M\), mit Ausnahme derjenigen, durch die eine Division bereits erfolglos versucht worden ist, mit \(\Omega\), so genügt es offenbar,
\SourceLeaf{408}
sämtliche Primteiler von \(M\), die in \(\Omega\) enthalten sind, zu finden. Weiß man nun auf irgendeinem Weg, dass eine Zahl \(r\), die kein Quadrat ist, quadratischer Rest von \(M\) ist, so kann keine Primzahl, für die \(r\) ein \(NR.\) ist, ein Teiler von \(M\) sein. Daher darf man aus \(\Omega\) sämtliche solchen Primzahlen entfernen, die meist ungefähr die Hälfte aller ausmachen. Weiß man außerdem von einer weiteren Zahl, die kein Quadrat ist, \(r'\), dass sie ein Rest von \(M\) ist, so kann man unter den nach dem ersten Ausschluss in \(\Omega\) verbleibenden Primzahlen wiederum diejenigen ausschließen, für die \(r'\) ein \(NR.\) ist. Diese machen wiederum ungefähr die Hälfte aus, sofern die Reste \(r\) und \(r'\) unabhängig sind, \textit{d. h.} sofern nicht der eine von selbst notwendig Rest sämtlicher Zahlen ist, von denen der andere Rest ist; dies wäre der Fall, wenn \(rr'\) ein Quadrat wäre. Sind noch weitere Reste von \(M\) bekannt, \(r'',\allowbreak r'''\) usw., die alle von den übrigen unabhängig sind\HistoricalNote{*)}{Ist das Produkt beliebig vieler der Zahlen \(r,\allowbreak r',\allowbreak r''\) usw. ein Quadrat, so ist jede von ihnen, etwa \(r\), Rest jeder Primzahl, die keine dieser Zahlen teilt und die Rest der übrigen Zahlen \(r',\allowbreak r''\) usw. ist. Damit also beliebig viele Reste als unabhängig betrachtet werden können, darf kein Produkt aus zweien, dreien usw. ein Quadrat sein.}, so lassen sich mit jedem ähnliche Ausschlüsse vornehmen. Dadurch nimmt die Menge der Zahlen in \(\Omega\) sehr rasch ab, sodass bald entweder alle entfernt sind, in welchem Fall \(M\) gewiss eine Primzahl ist, oder nur so wenige übrig bleiben, darunter offenbar sämtliche Primteiler von \(M\), falls es solche gibt, dass die Division durch sie ohne Mühe versucht werden kann. Bei einer Zahl, die eine Million nicht wesentlich übersteigt, genügen meist sechs oder sieben Ausschlüsse reichlich; bei einer Zahl mit acht oder neun Ziffern neun oder zehn. Über zwei Punkte müssen wir nun noch sprechen: \textit{erstens}, wie sich geeignete und genügend viele Reste von \(M\) finden lassen, und \textit{zweitens}, wie man den Ausschluss selbst am bequemsten ausführt. Wir werden jedoch die Reihenfolge dieser Fragen umkehren, zumal die zweite zeigen wird, welche Reste sich für diesen Zweck besonders eignen.
